\documentclass[12pt]{article}
\usepackage[utf8]{inputenc}
\usepackage[T1]{fontenc}
\usepackage[margin=1in]{geometry}
\usepackage{setspace}
\usepackage{xcolor}
\usepackage[colorlinks=true, urlcolor=blue!60!black, linkcolor=blue!60!black]{hyperref}
\setlength{\parindent}{0em}
\setlength{\parskip}{0.6em}
\renewcommand{\baselinestretch}{1.0}
\begin{document}
\begin{center}
{\LARGE \textit{Chiles v. Salazar}, 607 U.S. \_\_\_ (2026)\par}
\vspace{0.5em}
{\large Case Brief --- Two-Page Edition\par}
\end{center}

\section{Memory Jogger}
Colorado's ban on conversion therapy for minors, applied to a counselor's pure talk therapy, is viewpoint discrimination subject to strict scrutiny, not a regulation of professional conduct (\href{https://www.supremecourt.gov/opinions/25pdf/24-539_fd9g.pdf}{Slip Op.}).

\section{Detailed Case Facts}
Kaley Chiles, a licensed counselor, offers only talk therapy; some minor clients seek help to reduce unwanted attractions or feel at harmony with their bodies. Colorado's 2019 Minor Conversion Therapy Law bars efforts to change a minor's orientation or gender identity but permits ``[a]cceptance, support, and understanding'' and help with transition. Chiles challenged it only as applied to her talk therapy.

\section{Procedural History}
The district court found standing but denied a preliminary injunction under rational-basis review; the Tenth Circuit affirmed 2--1, with Judge Hartz calling the conduct label ``a labeling game.''

\section{Judicial Votes}
Reversed and remanded, 8--1. Gorsuch for the Court (Roberts, Thomas, Alito, Sotomayor, Kagan, Kavanaugh, Barrett); Kagan, with Sotomayor, concurring; Jackson dissenting.

\section{Holding}
As applied to pure talk therapy, the law regulates speech by viewpoint and triggers strict scrutiny. Relabeling speech as ``treatment'' or invoking a tradition of regulating ``substandard care'' does not avoid it. The Court remanded for application of strict scrutiny.

\section{Analysis of Opinions}
\textbf{Majority:} ``[A]ll Ms. Chiles does is speak.'' The law permits affirmation but forbids helping a minor pursue change. The speech-incident-to-conduct exception fails because others may give the same advice, and no historic tradition supports the claimed exemption; professional consensus cannot become enforced orthodoxy. \textbf{Kagan:} the law is ``textbook'' viewpoint discrimination, but viewpoint-neutral healthcare rules remain open. \textbf{Jackson:} States regulate clinical treatment, not ``speech as speech''; in medicine ``we insist upon competence, not debate.''

\section{Examples: Future Applications}
\textbf{Same-Side:} A mirror-image ban on affirming talk therapy faces strict scrutiny. \textbf{Opposite-Side:} A ban on electric shocks or nausea-inducing drugs regulates conduct and survives rational-basis review. \textbf{Fence-Sitter:} A viewpoint-neutral ban on steering minors toward any predetermined outcome divides the Kagan concurrence from the Jackson dissent.

\section{Critique}
Health-law critics call the speech-versus-procedure line clinically incoherent, since much of medicine is spoken advice. Free-speech advocates praise the decision for denying States power to enforce orthodoxy through ``medical consensus.''

\section{Key Quotations}
\begin{quote}
``Her speech does not become conduct just because the State may call it that. Nor does her speech become conduct just because it can also be described as a `treatment,' a `therapeutic modality,' or anything else. The First Amendment is no word game. And the rights it protects cannot be renamed away or their protections nullified by `mere labels.' ''
\end{quote}
\noindent --- Justice Gorsuch, Opinion of the Court, slip op. at 12--13.

\end{document}
